%%%%%%%%%%%%%%%%%%%%%%%%%%%%%%%%%%%%%%%%%%%%%%%%%%%%%%%%%%%%%%%%%%%%%%%%
%                                                                      %
% GLE - Graphics Layout Engine <http://glx.sourceforge.io/>            %
%                                                                      %
% Modified BSD License                                                 %
%                                                                      %
% Copyright (C) 2009 GLE.                                              %
%                                                                      %
% Redistribution and use in source and binary forms, with or without   %
% modification, are permitted provided that the following conditions   %
% are met:                                                             %
%                                                                      %
%    1. Redistributions of source code must retain the above copyright %
% notice, this list of conditions and the following disclaimer.        %
%                                                                      %
%    2. Redistributions in binary form must reproduce the above        %
% copyright notice, this list of conditions and the following          %
% disclaimer in the documentation and/or other materials provided with %
% the distribution.                                                    %
%                                                                      %
%    3. The name of the author may not be used to endorse or promote   %
% products derived from this software without specific prior written   %
% permission.                                                          %
%                                                                      %
% THIS SOFTWARE IS PROVIDED BY THE AUTHOR "AS IS" AND ANY EXPRESS OR   %
% IMPLIED WARRANTIES, INCLUDING, BUT NOT LIMITED TO, THE IMPLIED       %
% WARRANTIES OF MERCHANTABILITY AND FITNESS FOR A PARTICULAR PURPOSE   %
% ARE DISCLAIMED. IN NO EVENT SHALL THE AUTHOR BE LIABLE FOR ANY       %
% DIRECT, INDIRECT, INCIDENTAL, SPECIAL, EXEMPLARY, OR CONSEQUENTIAL   %
% DAMAGES (INCLUDING, BUT NOT LIMITED TO, PROCUREMENT OF SUBSTITUTE    %
% GOODS OR SERVICES; LOSS OF USE, DATA, OR PROFITS; OR BUSINESS        %
% INTERRUPTION) HOWEVER CAUSED AND ON ANY THEORY OF LIABILITY, WHETHER %
% IN CONTRACT, STRICT LIABILITY, OR TORT (INCLUDING NEGLIGENCE OR      %
% OTHERWISE) ARISING IN ANY WAY OUT OF THE USE OF THIS SOFTWARE, EVEN  %
% IF ADVISED OF THE POSSIBILITY OF SUCH DAMAGE.                        %
%                                                                      %
%%%%%%%%%%%%%%%%%%%%%%%%%%%%%%%%%%%%%%%%%%%%%%%%%%%%%%%%%%%%%%%%%%%%%%%%

\appendix
\chapter{Tables}
\section{Markers}
\index{markers}
\mbox{\input{primitives/fig/gc_marker.inc}}

\section{Functions and Variables}
\index{functions}
\label{fct:sec}

GLE has the following built in functions and variables in the following categories

\subsection{General Program}

\begin{supertabular}{lp{10cm}} \hline
Function name     & Returns  \\ \hline
{\tt arg(i)} \index{arg()}       	    & i-th command line argument \\
{\tt arg\$(i)} \index{arg()}            & i-th command line argument \\
{\tt date\$()} \index{date\$()}      	& current date e.g. ``Tue Apr 09 1991'' \\
{\tt device\$()} \index{device\$()}  	& available devices e.g. ``HARDCOPY, PS''\\
{\tt eval(str)} \index{eval()}         	& evaluates given GLE expression \\
{\tt exit [status]} \index{exit} & Terminates execution of the GLE script. The optional {\tt status} is an integer exit status returned by GLE to the operating system. If no status is specified, an exit status of zero is returned. \\
{\tt $\backslash{}$expr(exp)} \index{$\backslash{}$expr(exp)} & substitute result of evaluating ``exp'' \\
{\tt file\$()} \index{file\$()}         & returns GLE file name (without extension) \\
{\tt getenv(str)} \index{getenv{}}      & returns environment variable ``str''\\
{\tt height(name\$)} \index{height()} 	& the height of the object {\tt name\$} \\
{\tt is\_cairo()} \index{is\_cairo} & Returns {\tt TRUE} if Cairo is the currently selected rendering engine and {\tt FALSE} otherwise. \\
{\tt is\_ghostscript()} \index{is\_ghostscript} & Returns {\tt TRUE} if Ghostscript is the currently selected rendering engine and {\tt FALSE} otherwise. \\
{\tt nargs()} \index{nargs()}       	& number of command line arguments \\
{\tt pageheight()} \index{pageheight()} & the height of the page (from size command) \\
{\tt pagewidth()} \index{pagewidth()}	& the width of the page (from size command) \\
{\tt path\$()} \index{path\$()}         & returns the directory where the script is located \\
{\tt pointx(pt)} \index{pointx()}    	& the x value of named point pt \\
{\tt pointy(pt)} \index{pointy()}    	& the y value of named point pt \\
{\tt print expression...} \index{print} & Evaluates and prints one or more expressions to the console. Multiple expressions may be supplied on the same statement and are printed in order. \\
{\tt print\_color color expression...} \index{print\_color} & Prints one or more expressions to the console using the specified GLE color. The color is interpreted using GLE's normal color syntax and predefined color names. When console-color support is disabled at compile time, the command prints the text without ANSI color codes. \\
{\tt ptx(pt)} \index{ptx()}          	& the x value of named point pt (abbreviation for pointx) \\
{\tt pty(pt)} \index{pty()}         	& the y value of named point pt (abbreviation for pointy)\\
{\tt tdepth(str\$)} \index{tdepth()} 	& the depth of {\tt str\$} assuming current the font, size \\
{\tt theight(str\$)} \index{theight()}  & the height of {\tt str\$} assuming  current font, size \\
{\tt twidth(str\$)} \index{twidth()} 	& the width of {\tt str\$} assuming current font, size \\
{\tt time\$()} \index{time\$()}      	& current time e.g. ``11:44:27'' \\
{\tt width(name\$)} \index{width()} 	& the width of the object {\tt name\$} \\
{\tt xend()} \index{xend()}          	& the x end point of a text string when drawn \\
{\tt xpos()} \index{xpos()}          	& the current x point \\
{\tt yend()} \index{yend()}          	& the y end point of a text string when drawn \\
{\tt ypos()} \index{ypos()}          	& the current y point \\
\end{supertabular}


\subsection{Colors \& Color Space Conversions}

\begin{supertabular}{l p{10cm}}
\hline Function name & Returns \\
\hline
{\tt rgb(red,green,blue)} \index{rgb()} & create color given RGB values in the range 0 to 1 \\
{\tt rgba(red,green,blue,alpha)} \index{rgba()} & create color given RGB and alpha values in the range 0 to 1 \\
{\tt rgb255(red,green,blue)} \index{rgb255()} & create color given RGB values in the range 0 to 255 \\
{\tt rgba255(red,green,blue,alpha)} \index{rgba255()} & create color given RGB and alpha in the range 0 to 255 \\
{\tt hsv(hue,saturation,value)} \index{hsv()} & create color using HSV values in the range 0 to 1, except for hue which ranges from 0 to 360 degrees\\
{\tt hsl(hue,saturation,lightness)} \index{hsl()} & create color using HSL values in the range 0 to 1, except for hue which ranges from 0 to 360 degrees \\
{\tt cmyk(cyan,magenta,yellow,black)} \index{cmyk()} & create color given CMYK values in the range 0 to 1 \\
{\tt cvtgray($x$)} \index{cvtgray()} & create grayscale color, 0 = black and 1 = white \\
{\tt cvtrgb(red,green,blue)} \index{cvtrgb()} & same as rgb(red,green,blue) \\
{\tt cvtcolor(name)} \index{cvtcolor()} & convert a color name to its color value \\
{\tt colorscale(value,scale\$)} \index{colorscale()} & create color from a named color scale, where value ranges from 0 to 1 The supported color scales are {\tt GRAY}, {\tt GREY}, {\tt RED}, {\tt GREEN}, {\tt BLUE}, {\tt CYAN}, {\tt MAGENTA}, {\tt YELLOW}, {\tt VERDIGRIS}, {\tt THERMAL}, {\tt COOL}, {\tt RAINBOW}, {\tt VIRIDIS}, {\tt PLASMA}, {\tt INFERNO}, {\tt MAGMA}, {\tt BLUE\_WHITE\_RED}, and {\tt RED\_WHITE\_BLUE}. Values below 0 or above 1 are clamped to the color scale endpoints. \\
{\tt color\_intensity(color)} \index{color\_intensity()} & return the perceived intensity of a color on a scale from 0 to 255, calculated as $255(0.299R + 0.587G + 0.114B)$, where $R$, $G$, and $B$ are the normalized red, green, and blue components of the color \\
\end{supertabular}\\
\vspace{2ex}


Note: RGB, RGBA, HSV, HSL, and CMYK color components are normalized to the range 0 to 1, except for hue, which is specified in degrees from 0 to 360, and RGB255/RGBA255 components, which range from 0 to 255.

\subsection{String or Text Manipulation}

\begin{supertabular}{ll} \hline
Function Name     & Returns  \\ \hline
{\tt num\$(exp)} \index{num\$()}     	& string representation of {\tt exp} \\
{\tt num1\$(exp)} \index{num1\$()}   	& as above but with no spaces \\
{\tt pos(str1\$,str2\$,exp)} \index{pos()}     & position of {\tt str2\$} in {\tt str1\$} from {\tt exp} \\
{\tt right\$(str\$,exp)} \index{right\$()}     & rest of {\tt str\$} starting at {\tt exp} \\
{\tt seg\$(str\$,exp1,exp2)} \index{seg\$()}   & {\tt str\$} from {\tt exp1} to {\tt exp2} \\
{\tt val(str\$)} \index{val()}       	& value of the string {\tt str\$} \\
{\tt left\$(str\$,exp)} \index{left\$()}  & left {\tt exp} characters of {\tt str\$} \\
{\tt len(str\$)} \index{len()}       	& the length of {\tt str\$} \\
{\tt format\$(exp,format)} \index{format\$()} & format {\tt exp} as specified in {\tt format} (p.~\pageref{formatnum:pg})\\
\end{supertabular}


\subsection{Data Series} \label{sec:dataseries}

GLE provides functions for querying, copying, manipulating, and analyzing data series in a graph block. Data-series operations do not modify their input data series. Instead, functions that operate on an entire data series return a new, materialized data series. Consequently, data-series operations may be nested and chained. Input data series must be materialized before they can be used by these functions. Functions that return a data series are used on the right-hand side of a {\tt let} statement inside a graph block. The returned data series is materialized immediately; the {\tt materialize} option to {\tt let} is therefore not required. For example:

\begin{verbatim}
begin graph
    ...
    data "data.dat" d1=c1,c2
    let d2 = smooth(d1,5)
    let d3 = normalize(d2)
    let d4 = derivative(d3)
    ! let generated data series must use maerialize
    let d10 = sin(x) from 0 to 10*pi nsteps 1000 MATERIALIZE
    d11 = derivative(d10)
    ...
end graph
\end{verbatim}

Data-series operations may also be nested:

\begin{verbatim}
    let d2 = normalize(smooth(d1,5))
\end{verbatim}

The following functions are available:

\begin{supertabular}{lp{10cm}}
\hline Function Name & Returns \\
\hline {\tt ndata(ds\$)} \index{ndata()} & The number of data points in the data series named by {\tt ds\$}. Returns zero if the data series does not exist. \\
{\tt dataxvalue(ds\$,i)} \index{dataxvalue()} & The x-coordinate of point {\tt i} in the data series named by {\tt ds\$}. Data-point indices start at 1. An error is generated if {\tt i} is outside the range 1 through {\tt ndata(ds\$)}. \\
{\tt datayvalue(ds\$,i)} \index{datayvalue()} & The y-coordinate of point {\tt i} in the data series named by {\tt ds\$}. Data-point indices start at 1. An error is generated if {\tt i} is outside the range 1 through {\tt ndata(ds\$)}. \\
{\tt copy(ds)} \index{copy()} & A new data series containing a copy of {\tt ds}. The returned data series is independent of the source data series. \\
{\tt smooth(ds,n)} \index{smooth()} & A new data series in which the y-values of {\tt ds} are replaced by moving averages over a window of {\tt n} points. The x-values are preserved. {\tt n} must be a positive integer. The smoothing window is shortened at the ends of the data series and at gaps in the data. The source data series is unchanged. \\
{\tt normalize(ds)} \index{normalize()} & A new data series whose valid y-values are linearly normalized to the range 0 to 1. The minimum y-value is mapped to 0 and the maximum y-value is mapped to 1. If all valid y-values are equal, they are mapped to zero. The x-values and point ordering are preserved. \\
{\tt clip(ds,ymin,ymax)} \index{clip()} & A new data series whose y-values are restricted to the interval $[y_{\rm min},y_{\rm max}]$. Values below {\tt ymin} are replaced by {\tt ymin}, and values above {\tt ymax} are replaced by {\tt ymax}. The x-values and number of data points are unchanged. \\
{\tt slice(ds,first,last)} \index{slice()} & A new data series containing points {\tt first} through {\tt last} of {\tt ds}, inclusive. GLE data-series indices start at 1, so {\tt slice(ds,1,ndata("ds"))} selects the complete data series. The selected x- and y-values are preserved. \\
{\tt derivative(ds)} \index{derivative()} & A new data series containing a numerical estimate of $dy/dx$. The original x-values are preserved. Interior derivatives use a spacing-aware three-point estimate suitable for nonuniformly spaced x-values. One-sided differences are used at the ends of a valid run of data. Missing or invalid points separate the data into independent runs. \\
{\tt integrate(ds)} \index{integrate()} & A new data series containing the cumulative integral of {\tt ds}. Integration is performed using the trapezoidal rule and the original x-values are preserved. The cumulative integral begins at zero and restarts at zero after a gap in the data. \\
{\tt integral(ds,x1,x2)} \index{integral()} & The scalar definite integral of {\tt ds} from x-coordinate {\tt x1} to {\tt x2}. Integration uses the trapezoidal rule applied to the piecewise linear representation of the data. If an integration bound falls between two data points, the y-value at the bound is obtained by linear interpolation. Reversing the bounds reverses the sign of the result. The x-values must be finite and strictly increasing, and the integration interval must lie within the x-range of the data series. \\
\hline \end{supertabular}


\subsection{Logical Operators}

\begin{supertabular}{lp{10cm}}
\hline
Operator & Meaning  \\
 \hline
{\tt =} \index{=}      & equals (same as assignment) \\
{\tt \textless \textgreater} \index{\textless \textgreater}  & not equals (in GLE {\tt !} is the comment character)  \\
{\tt \textless } \index{\textless} & less than\\
{\tt \textgreater } \index{\textgreater}  & greater than\\
{\tt \textless =} \index{\textless =}  & less than or equal \\
{\tt \textgreater =} \index{\textgreater =}  & greater than or equal \\
{\tt and} \index{and}     & logical and between two expressions e.g. {\tt ( x \textless\ 1) and  (x \textgreater\ 10) }\\
{\tt or} \index{or}     & logical or between two expressions\\
{\tt not(exp)} \index{not()}   & logical not of {\tt exp} \\
\end{supertabular}

\subsection{Mathematical Operators}

GLE provides the following mathematical operators, functions, and constants.  All angles are in radians. The constant values have 16 digits of precision within GLE.  The special functions are from the boost C++ libraries or the C++ standard math libraries.  The physical constants are based on CODATA values and are in SI units and generated from scipy.constants.  The mathematical constants are from the boost::math C++ library.

\vspace{2ex}

\begin{tabular}{ll} \hline
Operator & Meaning  \\ \hline
{\tt =} \index{=}    & assignment (same as logical equals) \\
{\tt +} \index{+}    & addition \\
{\tt -} \index{-}    & subtraction \\
{\tt *} \index{*}    & multiplication\\
{\tt /} \index{/}    & division \\
{\tt \textasciicircum } \index{\textasciicircum}  & exponentiation \\
{\tt \%} \index{\%}  & modulus or remainder\\
{\tt ++} \index{++}  & increment {\tt a++ $\equiv$ a = a + 1 }\\
{\tt --} \index{--}  & decrement {\tt a-- $\equiv$ a = a - 1 }\\
{\tt +=} \index{+=}  & addition assignment  {\tt a+=3 $\equiv$ a = a + 3 }\\
{\tt -=} \index{-=}  & subtraction assignment {\tt a-=3 $\equiv$ a = a - 3 }\\
{\tt *=} \index{*=}  & multiplication assignment {\tt a*=3 $\equiv$ a = a * 3 }\\
{\tt /=} \index{/=}  & division assignment {\tt a/=3 $\equiv$ a = a / 3 }\\
\end{tabular}

\vspace{2ex}

\subsection{Mathematical Constants}

\begin{tabular}{ll} \hline
Mathematical Constants & Value  \\ \hline
{\tt pi} \index{pi}      & 3.14159265358979323846  \\
{\tt mc\_catalan} \index{mc\_catalan}  & 0.9159655941772190 \\
{\tt mc\_degree} \index{mc\_degree}  & 0.0174532925199433 \\
{\tt mc\_e} \index{mc\_e}  & 2.7182818284590451 \\
{\tt mc\_e\_pow\_pi} \index{mc\_e\_pow\_pi}  & 23.1406926327792704 \\
{\tt mc\_euler} \index{mc\_euler}  & 0.5772156649015329 \\
{\tt mc\_exp\_minus\_half} \index{mc\_exp\_minus\_half}  & 0.6065306597126334 \\
{\tt mc\_half} \index{mc\_half}  & 0.5000000000000000 \\
{\tt mc\_half\_pi} \index{mc\_half\_pi}  & 1.5707963267948966 \\
{\tt mc\_half\_root\_two} \index{mc\_half\_root\_two}  & 0.7071067811865476 \\
{\tt mc\_ln\_phi} \index{mc\_ln\_phi}  & 0.4812118250596035 \\
{\tt mc\_one\_div\_two\_pi} \index{mc\_one\_div\_two\_pi}  & 0.1591549430918953 \\
{\tt mc\_half\_root\_two} \index{mc\_half\_root\_two}  & 0.7071067811865476 \\
{\tt mc\_phi} \index{mc\_phi}  & 1.6180339887498949 \\
{\tt mc\_pi} \index{mc\_pi}  & 3.1415926535897931 \\
{\tt mc\_quarter\_pi} \index{mc\_quarter\_pi}  & 0.7853981633974483 \\
{\tt mc\_radian} \index{mc\_radian}  & 57.2957795130823229 \\
{\tt mc\_root\_e} \index{mc\_root\_e}  & 1.6487212707001282 \\
{\tt mc\_root\_half\_pi} \index{mc\_root\_half\_pi}  & 1.2533141373155003 \\
{\tt mc\_root\_ln\_four} \index{mc\_root\_ln\_four}  & 1.1774100225154747 \\
{\tt mc\_root\_pi} \index{mc\_root\_pi}  & 1.7724538509055161 \\
{\tt mc\_root\_three} \index{mc\_root\_three}  & 1.7320508075688772 \\
{\tt mc\_root\_two} \index{mc\_root\_two}  & 1.4142135623730951 \\
{\tt mc\_root\_two\_pi} \index{mc\_root\_two\_pi}  & 2.5066282746310007 \\
{\tt mc\_third} \index{mc\_third}  & 0.3333333333333333 \\
{\tt mc\_third\_pi} \index{mc\_third\_pi}  & 1.0471975511965979 \\
{\tt mc\_two\_pi} \index{mc\_two\_pi}  & 6.2831853071795862 \\
{\tt mc\_two\_thirds} \index{mc\_two\_thirds}  & 0.6666666666666666 \\
{\tt mc\_zeta\_three} \index{mc\_zeta\_three}  & 1.2020569031595942 \\
{\tt mc\_zeta\_two} \index{mc\_zeta\_two}  & 1.6449340668482264 \\
\end{tabular}


\begin{tabular}{lll} \hline
Physical Constants & Value & Units \\ \hline
{\tt pc\_c} \index{pc\_c}  & 299792458.0 & SI \\
{\tt pc\_speed\_of\_light} \index{pc\_speed\_of\_light}  & 299792458.0 & SI \\
{\tt pc\_mu\_0} \index{pc\_mu\_0}  & 1.25663706127e-06 & SI \\
{\tt pc\_epsilon\_0} \index{pc\_epsilon\_0}  & 8.8541878188e-12 & SI \\
{\tt pc\_h} \index{pc\_h}  & 6.62607015e-34 & SI \\
{\tt pc\_Planck} \index{pc\_Planck}  & 6.62607015e-34 & SI \\
{\tt pc\_hbar} \index{pc\_hbar}  & 1.0545718176461565e-34 & SI \\
{\tt pc\_gravitational\_constant} \index{pc\_gravitational\_constant}  & 6.6743e-11 & SI \\
{\tt pc\_g} \index{pc\_g}  & 9.80665 & SI \\
{\tt pc\_e} \index{pc\_e}  & 1.602176634e-19 & SI \\
{\tt pc\_elementary\_charge} \index{pc\_elementary\_charge}  & 1.602176634e-19 & SI \\
{\tt pc\_R} \index{pc\_R}  & 8.31446261815324 & SI \\
{\tt pc\_gas\_constant} \index{pc\_gas\_constant}  & 8.31446261815324 & SI \\
{\tt pc\_alpha} \index{pc\_alpha}  & 0.0072973525643 & SI \\
{\tt pc\_fine\_structure} \index{pc\_fine\_structure}  & 0.0072973525643 & SI \\
{\tt pc\_N\_A} \index{pc\_N\_A}  & 6.02214076e+23 & SI \\
{\tt pc\_Avogadro} \index{pc\_Avogadro}  & 6.02214076e+23 & SI \\
{\tt pc\_k} \index{pc\_k}  & 1.380649e-23 & SI \\
{\tt pc\_Boltzmann} \index{pc\_Boltzmann}  & 1.380649e-23 & SI \\
{\tt pc\_sigma} \index{pc\_sigma}  & 5.6703744191844314e-08 & SI \\
{\tt pc\_Stefan\_Boltzmann} \index{pc\_Stefan\_Boltzmann}  & 5.6703744191844314e-08 & SI \\
{\tt pc\_Wien} \index{pc\_Wien}  & 0.0028977719551851727 & SI \\
{\tt pc\_Rydberg} \index{pc\_Rydberg}  & 10973731.568157 & SI \\
{\tt pc\_m\_e} \index{pc\_m\_e}  & 9.1093837139e-31 & SI \\
{\tt pc\_electron\_mass} \index{pc\_electron\_mass}  & 9.1093837139e-31 & SI \\
{\tt pc\_m\_p} \index{pc\_m\_p}  & 1.67262192595e-27 & SI \\
{\tt pc\_proton\_mass} \index{pc\_proton\_mass}  & 1.67262192595e-27 & SI \\
{\tt pc\_m\_n} \index{pc\_m\_n}  & 1.67492750056e-27 & SI \\
{\tt pc\_neutron\_mass} \index{pc\_neutron\_mass}  & 1.67492750056e-27 & SI \\
{\tt pc\_zero\_Celsius} \index{pc\_zero\_Celsius}  & 273.15 & SI \\
\end{tabular}


\subsection{CODATA Physical Constants} \index{CODATA} \index{physical constants}

GLE provides access to the physical constants in the CODATA database through the functions \texttt{codata()}, \texttt{codata\_units()}, and \texttt{codata\_uncertainty()}. The constants are identified by their CODATA names.

\begin{verbatim}
codata(name$) codata_units(name$) codata_uncertainty(name$)
\end{verbatim}
\begin{description}
\item[\texttt{codata(name\$)}] Returns the numerical value of the physical constant specified by \texttt{name\$}.
\item[\texttt{codata\_units(name\$)}] Returns a string containing the units of the physical constant specified by \texttt{name\$}.
\item[\texttt{codata\_uncertainty(name\$)}] Returns the standard uncertainty of the physical constant specified by \texttt{name\$}.
\end{description}

For example:
\begin{verbatim}
c = codata("speed of light in vacuum")
u$ = codata_units("speed of light in vacuum")
print c u$
\end{verbatim}

prints the value of the speed of light and its corresponding units:

\begin{verbatim}
299792458 m s^-1
\end{verbatim}

The uncertainty of a constant can be obtained similarly:

\begin{verbatim}
m = codata("electron mass")
mu = codata_uncertainty("electron mass")
print m mu
\end{verbatim}

The names used by these functions correspond to the names in the CODATA physical constants database. Consult SciPy docuementation for list of allconstants. For example:

\begin{verbatim}
 codata("Planck constant")
 codata("electron mass")
 codata("proton mass")
 codata("Boltzmann constant")
 codata("speed of light in vacuum")
\end{verbatim}

The CODATA data are incorporated into GLE when the source code is generated. GLE therefore does not require SciPy or any external CODATA database at runtime. The existing \texttt{PC\_} constants provide convenient names for a selected set of commonly used physical constants. The \texttt{codata()} functions provide access to the complete CODATA physical constants database, including the associated units and uncertainties.

\vspace{2ex}

\subsection{Mathematical Functions}

\begin{supertabular}{lp{10cm}}
\hline Function & Returns \\
 \hline
  {\tt abs($x$)} \index{abs()} & absolute value of $x$ \\
  {\tt acos($x$)} \index{acos()} & inverse cosine of $x$ \\
  {\tt acosh($x$)} \index{acosh()} & inverse hyperbolic cosine of $x$ \\
  {\tt acot($x$)} \index{acot()} & 1/atan($x$) \\
  {\tt acoth($x$)} \index{acoth()} & 1/atanh($x$) \\
  {\tt acsc($x$)} \index{acsc()} & 1/asin($x$) \\
  {\tt acsch($x$)} \index{acsch()} & 1/asinh($x$) \\
  {\tt asec($x$)} \index{asec()} & 1/acos($x$) \\
  {\tt asech($x$)} \index{asech()} & 1/acosh($x$) \\
  {\tt asin($x$)} \index{asin()} & inverse sine of $x$ \\
  {\tt asinh($x$)} \index{asinh()} & inverse hyperbolic sine of $x$ \\
  {\tt atan($x$)} \index{atan()} & inverse tangent of $x$ \\
  {\tt atanh($x$)} \index{atanh()} & inverse hyperbolic tangent of $x$ \\
  {\tt atn($x$)} \index{atn()} & same as atan($x$) ({\it deprecated, kept for backward compatibility}) \\
  {\tt cbrt($x$)} \index{cbrt()} & cube root of $x$ \\
  {\tt ceil($x$)} \index{ceil()} & smallest integer greater than or equal to $x$ \\
  {\tt clamp(value,min,max)} \index{clamp()} & return value constrained to the range min to max. An error is   generated if min is greater than max \\
  {\tt cos($x$)} \index{cos()} & cosine of $x$ \\
  {\tt cosh($x$)} \index{cosh()} & hyperbolic cosine of $x$ \\
  {\tt cot($x$)} \index{cot()} & 1/tan($x$) \\
  {\tt coth($x$)} \index{coth()} & 1/tanh($x$) \\
  {\tt csc($x$)} \index{csc()} & 1/sin($x$) \\
  {\tt csch($x$)} \index{csch()} & 1/sinh($x$) \\
  {\tt erf($x$)} \index{erf()} & Gaussian error function of $x$ \\
  {\tt erfc($x$)} \index{erfc()} & complementary Gaussian error function of $x$ \\
  {\tt exp($x$)} \index{exp()} & $e^x$ \\
  {\tt expm1($x$)} \index{expm1()} & $e^x-1$, computed accurately for $x$ near zero \\
  {\tt fix($x$)} \index{fix()} & $x$ rounded towards 0 \\
  {\tt floor($x$)} \index{floor()} & largest integer less than or equal to $x$ \\
  {\tt gamma($x$)} \index{gamma()} & Gamma function $\Gamma(x)$ \\
  {\tt hypot($x$,$y$)} \index{hypot()} & $\sqrt{x^2+y^2}$, computed to avoid undue overflow and underflow \\
  {\tt int($x$)} \index{int()} & integer part of $x$ \\
  {\tt is\_close(a,b,rel\_tol)} \index{is\_close()} & return true if a and b are equal within the relative tolerance rel\_tol, where the comparison is $|a-b| \leq {\tt rel\_tol}\max(|a|,|b|)$ \\
  {\tt is\_close\_abs(a,b,abs\_tol)} \index{is\_close\_abs()} & return true if a and b are equal within the absolute tolerance abs\_tol, where the comparison is $|a-b| \leq {\tt abs\_tol}$ \\
  {\tt lerp(a,b,t)} \index{lerp()} & return the linear interpolation between $a$ and $b$ at $t$, calculated as $a+t(b-a)$. Values of $t$ outside the range 0 to 1 perform linear extrapolation \\
  {\tt log($x$)} \index{log()} & log to base $e$ of $x$ \\
  {\tt log10($x$)} \index{log10()} & log to base 10 of $x$ \\
  {\tt log1p($x$)} \index{log1p()} & $\log(1+x)$, computed accurately for $x$ near zero \\
  {\tt log\_gamma($x$)} \index{log\_gamma()} & natural logarithm of the absolute value of $\Gamma(x)$ \\
  {\tt max($x$,$y$)} \index{max()} & the maximum value of $x$ or $y$ \\
  {\tt min($x$,$y$)} \index{min()} & the minimum value of $x$ or $y$ \\
  {\tt random($x$)} \index{random()} & random number between 0 and $x$ \\
  {\tt linear\_interpolate($x$,$x_1$,$x_2$,$y_1$,$y_2$)} \index{linear\_interpolate()} & Interpolates the value $x$ from the input range $x_1$ to $x_2$ onto the output range $y_1$ to $y_2$. Values outside the input range are extrapolated, and reversed ranges are supported. An error is generated if $x_1$ equals $x_2$. \\
  {\tt rnd($x$)} \index{rnd()} & same as random($x$) ({\it retained for backward compatibility}) \\
  {\tt round($x$)} \index{round()} & $x$ rounded to the nearest integer, halfway cases away from zero \\
  {\tt sdiv($x$,$y$)} \index{sdiv()} & return $x/y$ or 0 if $y=0$ \\
  {\tt sec($x$)} \index{sec()} & 1/cos($x$) \\
  {\tt sech($x$)} \index{sech()} & 1/cosh($x$) \\
  {\tt sgn($x$)} \index{sgn()} & returns 1 if $x$ is positive, -1 if $x$ is negative \\
  {\tt sin($x$)} \index{sin()} & sine of $x$ \\
  {\tt sinc($x$)} \index{sinc()} & sinus cardinal function, $\sin(x)/x$, with sinc(0) = 1 \\
  {\tt sinh($x$)} \index{sinh()} & hyperbolic sine of $x$ \\
  {\tt sqr($x$)} \index{sqr()} & $x$ squared \\
  {\tt sqrt($x$)} \index{sqrt()} & square root of $x$ \\
  {\tt tan($x$)} \index{tan()} & tangent of $x$ \\
  {\tt tanh($x$)} \index{tanh()} & hyperbolic tangent of $x$ \\
  {\tt todeg($x$)} \index{todeg()} & convert from radians to degrees \\
  {\tt torad($x$)} \index{torad()} & convert from degrees to radians \\
  {\tt airy\_first($x$)} \index{airy\_first()} & Airy function of the first kind \\
  {\tt airy\_second($x$)} \index{airy\_second()} & Airy function of the second kind \\
  {\tt associated\_laguerre($n$,$x$)} \index{associated\_laguerre()} & $n^{th}$ order associated Laguerre polynomial of $x$ \\
  {\tt associated\_legendre($\ell$,$x$)} \index{associated\_legendre()} & $\ell^{th}$-order associated Legendre polynomial of $x$ \\
  {\tt bessel\_first($v$,$x$)} \index{bessel\_first()} & Bessel function of the first kind $J_v(x)$ \\
  {\tt bessel\_modified\_first($v$,$x$)} \index{bessel\_modified\_first()} & modified Bessel function of the first kind $I_v(x)$ \\
  {\tt bessel\_modified\_second($v$,$x$)} \index{bessel\_modified\_second()} & modified Bessel function of the second kind $K_v(x)$ \\
  {\tt bessel\_second($v$,$x$)} \index{bessel\_second()} & Bessel function of the second kind $Y_v(x)$ \\
  {\tt chebyshev\_first($n$,$x$)} \index{chebyshev\_first()} & $n^{th}$ order Chebyshev polynomial of the first kind of $x$ \\
  {\tt chebyshev\_second($n$,$x$)} \index{chebyshev\_second()} & $n^{th}$ order Chebyshev polynomial of the second kind of $x$ \\
  {\tt double\_factorial($n$)} \index{double\_factorial()} & $n!!$ \\
  {\tt factorial($n$)} \index{factorial()} & $n!$ \\
  {\tt hermite($n$,$x$)} \index{hermite()} & $n^{th}$ order Hermite polynomial of $x$ \\
  {\tt sph\_bessel($n$,$x$)} \index{sph\_bessel()} & same as spherical\_bessel\_first($n$,$x$) \\
  {\tt sph\_neumann($n$,$x$)} \index{sph\_neumann()} & same as spherical\_bessel\_second($n$,$x$) \\
  {\tt spherical\_bessel\_first($n$,$x$)} \index{spherical\_bessel\_first()} & spherical Bessel function of the first kind $j_n(x)$ \\
  {\tt spherical\_bessel\_second($n$,$x$)} \index{spherical\_bessel\_second()} & spherical Bessel function of the second kind $y_n(x)$ \\
  {\tt spherical\_harmonic($n$,$m$,$\theta$,$\phi$)} \index{spherical\_harmonic()} & real valued spherical harmonic, $\theta$ polar angle $[0,\pi]$, $\phi$ azimuthal angle $[0,2\pi]$ \\
\end{supertabular}

\subsection{Graphing}

Graph variables and function that provide information from preceding graph block

\begin{supertabular}{lp{10cm}} \hline
Function Name     & Returns  \\ \hline
{\tt xgmin} \index{xgmin} & the minimum x-coordinate of the graph \\
{\tt xgmax} \index{xgmax} & the maximum x-coordinate of the graph \\
{\tt x2gmin} \index{x2gmin} & the minimum x2-coordinate of the graph \\
{\tt x2gmax} \index{x2gmax} & the maximum x2-coordinate of the graph \\
{\tt ygmin} \index{ygmin} & the minimum y-coordinate of the graph \\
{\tt ygmax} \index{ygmax} & the maximum y-coordinate of the graph \\
{\tt y2gmin} \index{y2gmin} & the minimum y2-coordinate of the graph \\
{\tt y2gmax} \index{y2gmax} & the maximum y2-coordinate of the graph \\
{\tt xbar(x,i)} \index{xbar()}        	& the absolute x coordinate of the i-th bar at point x on the graph \\
{\tt xg(xexp)} \index{xg()}             & converts units of last graph to abs cm. \\
{\tt xg3d(x,y,z)} \index{xg3d()}        & converts units of last 3D graph to abs cm. \\
{\tt xy2angle(dx,dy)} \index{xy2angle()}& convert rectangular coordinates to polar angle (in degrees)\\
{\tt yg(yexp)} \index{yg()}             & converts units of last graph to abs cm. \\
{\tt yg3d(x,y,z)} \index{yg3d()}        & converts units of last 3D graph to abs cm. \\

\end{supertabular}

%%%%%%%%%%%%%%%%%%%%%%%%%%%%%%%%%%%%%%%%%%%%%%%%%%%%%%%%%%%%%%%%%%%%%%%%
%                                                                      %
% GLE - Graphics Layout Engine <http://glx.sourceforge.io/>            %
%                                                                      %
% Modified BSD License                                                 %
%                                                                      %
% Copyright (C) 2009 GLE.                                              %
%                                                                      %
% Redistribution and use in source and binary forms, with or without   %
% modification, are permitted provided that the following conditions   %
% are met:                                                             %
%                                                                      %
%    1. Redistributions of source code must retain the above copyright %
% notice, this list of conditions and the following disclaimer.        %
%                                                                      %
%    2. Redistributions in binary form must reproduce the above        %
% copyright notice, this list of conditions and the following          %
% disclaimer in the documentation and/or other materials provided with %
% the distribution.                                                    %
%                                                                      %
%    3. The name of the author may not be used to endorse or promote   %
% products derived from this software without specific prior written   %
% permission.                                                          %
%                                                                      %
% THIS SOFTWARE IS PROVIDED BY THE AUTHOR "AS IS" AND ANY EXPRESS OR   %
% IMPLIED WARRANTIES, INCLUDING, BUT NOT LIMITED TO, THE IMPLIED       %
% WARRANTIES OF MERCHANTABILITY AND FITNESS FOR A PARTICULAR PURPOSE   %
% ARE DISCLAIMED. IN NO EVENT SHALL THE AUTHOR BE LIABLE FOR ANY       %
% DIRECT, INDIRECT, INCIDENTAL, SPECIAL, EXEMPLARY, OR CONSEQUENTIAL   %
% DAMAGES (INCLUDING, BUT NOT LIMITED TO, PROCUREMENT OF SUBSTITUTE    %
% GOODS OR SERVICES; LOSS OF USE, DATA, OR PROFITS; OR BUSINESS        %
% INTERRUPTION) HOWEVER CAUSED AND ON ANY THEORY OF LIABILITY, WHETHER %
% IN CONTRACT, STRICT LIABILITY, OR TORT (INCLUDING NEGLIGENCE OR      %
% OTHERWISE) ARISING IN ANY WAY OUT OF THE USE OF THIS SOFTWARE, EVEN  %
% IF ADVISED OF THE POSSIBILITY OF SUCH DAMAGE.                        %
%                                                                      %
%%%%%%%%%%%%%%%%%%%%%%%%%%%%%%%%%%%%%%%%%%%%%%%%%%%%%%%%%%%%%%%%%%%%%%%%

\newpage
\section{\LaTeX\ Macros}
\label{ltxsym:sec}
\index{\LaTeX!macros}

\index{\LaTeX} There are several \LaTeX\ like commands 
which can be used within text, they are:

\vspace{0.5cm}
\begin{tabular}{ll}
\verb+\` \' \v \u \= \^+    & Implemented TeX accents \\
\verb+\. \H \~ \''+         & Implemented TeX accents \\
\verb+^{}+                  & Superscript \\
\verb+_{}+                  & Subscript \\
\verb+\\+                   & Forced newline \\
\verb+\_+                   & Underscore character \\
\verb+\,+                   & 0.5em (em = width of the letter `m') \\
\verb+\:+                   & 1em space \\
\verb+\;+                   & 2em space \\
\verb+\!+                   & -0.5em space \\
\verb+\tex{expression}+     & Any LaTeX expression \\
\verb+\char{22}+            & Any character in current font \\
\verb+\chardef{a}{hello}+   & Define a character as a macro \\
\verb+\def\v{hello}+        & Defines a macro \\
\verb+\movexy{2}{3}+        & Moves the current text point \\
\verb+\glass+               & Makes move/space work on beginning of line \\
\verb+\it+                  & Switches to italic font \\ 
\verb+\kern{-0.1em}+        & Change inter character distance \\
\verb+\ldots+               & $\ldots$ \\
\verb+\lineskip{.1}+        & Sets the default distance between lines of text \\
\verb+\linegap{-1}+         & Sets the minimum required gap between lines \\
\verb+\lower{0.1em}{hello}+ & Lower the given text \\
\verb+\parskip{0.1em}+      & Set distance between paragraphs \\
\verb+\raise{0.1em}{hello}+ & Raise the given text \\
\verb+\rm+                  & Switches to roman font \\
\verb+\bf+                  & Switches to bold font \\
\verb+\rule{2}{4}+          & Draws a filled in box, 2cm by 4cm \\
\verb+\setfont{rmb}+        & Sets the current text font \\
\verb+\sethei{.3}+          & Sets the font height (in cm) \\
\verb+\setstretch{2}+       & Scales the quantity of glue between words \\
\verb+\tt+                  & Switches to typewriter (fixed space) font \\
\end{tabular}

\newpage
\section{\LaTeX\ Symbols}
\index{symbols}

\mbox{\input{appendix/fig/symbols.inc}}

\newpage

%\section{Installing GLE}
%\label{install:app}

\section{Fonts}

\index{fonts} \index{font-examples}
\resizebox{!}{22cm}{\includegraphics{appendix/fig/ap_fonts}}
\newpage

%%
% \clearpage
% %%%%%%%%%%%%%%%%%%%%%%%%%%%%%%%%%%%%%%%%%%%%%%%%%%%%%%%%%%%%%%%%%%%%%%%%
%                                                                      %
% GLE - Graphics Layout Engine <http://glx.sourceforge.io/>            %
%                                                                      %
% Modified BSD License                                                 %
%                                                                      %
% Copyright (C) 2009 GLE.                                              %
%                                                                      %
% Redistribution and use in source and binary forms, with or without   %
% modification, are permitted provided that the following conditions   %
% are met:                                                             %
%                                                                      %
%    1. Redistributions of source code must retain the above copyright %
% notice, this list of conditions and the following disclaimer.        %
%                                                                      %
%    2. Redistributions in binary form must reproduce the above        %
% copyright notice, this list of conditions and the following          %
% disclaimer in the documentation and/or other materials provided with %
% the distribution.                                                    %
%                                                                      %
%    3. The name of the author may not be used to endorse or promote   %
% products derived from this software without specific prior written   %
% permission.                                                          %
%                                                                      %
% THIS SOFTWARE IS PROVIDED BY THE AUTHOR "AS IS" AND ANY EXPRESS OR   %
% IMPLIED WARRANTIES, INCLUDING, BUT NOT LIMITED TO, THE IMPLIED       %
% WARRANTIES OF MERCHANTABILITY AND FITNESS FOR A PARTICULAR PURPOSE   %
% ARE DISCLAIMED. IN NO EVENT SHALL THE AUTHOR BE LIABLE FOR ANY       %
% DIRECT, INDIRECT, INCIDENTAL, SPECIAL, EXEMPLARY, OR CONSEQUENTIAL   %
% DAMAGES (INCLUDING, BUT NOT LIMITED TO, PROCUREMENT OF SUBSTITUTE    %
% GOODS OR SERVICES; LOSS OF USE, DATA, OR PROFITS; OR BUSINESS        %
% INTERRUPTION) HOWEVER CAUSED AND ON ANY THEORY OF LIABILITY, WHETHER %
% IN CONTRACT, STRICT LIABILITY, OR TORT (INCLUDING NEGLIGENCE OR      %
% OTHERWISE) ARISING IN ANY WAY OUT OF THE USE OF THIS SOFTWARE, EVEN  %
% IF ADVISED OF THE POSSIBILITY OF SUCH DAMAGE.                        %
%                                                                      %
%%%%%%%%%%%%%%%%%%%%%%%%%%%%%%%%%%%%%%%%%%%%%%%%%%%%%%%%%%%%%%%%%%%%%%%%

\addcontentsline{toc}{section}{A Single Graph}
\psgraphin{appendix/fig/single}{12.0}{9.0}{Single Graph}
%%\graphin{single}{12.0cm}{9.0cm}{1}

\vfill
%%\input{single.gle}
\begin{center}
\begin{minipage}[t]{9.0cm}
\begin{verbatim}
size 12 9  box
begin graph
	SIZE 12 9
	Title "Bar example"
	data single.dat
	bar d1,d2 fill green,blue 
	fill x1,d1 color grey10
end graph
\end{verbatim}
\hfill {\small single.gle}\\
\vfill
%%\input{single.dat}
\begin{verbatim}
0 2 6
1 10 16
2 3 6 
3 6 3
4 5 4 
5 3 7 
\end{verbatim}
\hfill {\small single.dat}\\
\end{minipage}
\end{center}
%%\vfill

%%
% \clearpage
% %%%%%%%%%%%%%%%%%%%%%%%%%%%%%%%%%%%%%%%%%%%%%%%%%%%%%%%%%%%%%%%%%%%%%%%%
%                                                                      %
% GLE - Graphics Layout Engine <http://glx.sourceforge.io/>            %
%                                                                      %
% Modified BSD License                                                 %
%                                                                      %
% Copyright (C) 2009 GLE.                                              %
%                                                                      %
% Redistribution and use in source and binary forms, with or without   %
% modification, are permitted provided that the following conditions   %
% are met:                                                             %
%                                                                      %
%    1. Redistributions of source code must retain the above copyright %
% notice, this list of conditions and the following disclaimer.        %
%                                                                      %
%    2. Redistributions in binary form must reproduce the above        %
% copyright notice, this list of conditions and the following          %
% disclaimer in the documentation and/or other materials provided with %
% the distribution.                                                    %
%                                                                      %
%    3. The name of the author may not be used to endorse or promote   %
% products derived from this software without specific prior written   %
% permission.                                                          %
%                                                                      %
% THIS SOFTWARE IS PROVIDED BY THE AUTHOR "AS IS" AND ANY EXPRESS OR   %
% IMPLIED WARRANTIES, INCLUDING, BUT NOT LIMITED TO, THE IMPLIED       %
% WARRANTIES OF MERCHANTABILITY AND FITNESS FOR A PARTICULAR PURPOSE   %
% ARE DISCLAIMED. IN NO EVENT SHALL THE AUTHOR BE LIABLE FOR ANY       %
% DIRECT, INDIRECT, INCIDENTAL, SPECIAL, EXEMPLARY, OR CONSEQUENTIAL   %
% DAMAGES (INCLUDING, BUT NOT LIMITED TO, PROCUREMENT OF SUBSTITUTE    %
% GOODS OR SERVICES; LOSS OF USE, DATA, OR PROFITS; OR BUSINESS        %
% INTERRUPTION) HOWEVER CAUSED AND ON ANY THEORY OF LIABILITY, WHETHER %
% IN CONTRACT, STRICT LIABILITY, OR TORT (INCLUDING NEGLIGENCE OR      %
% OTHERWISE) ARISING IN ANY WAY OUT OF THE USE OF THIS SOFTWARE, EVEN  %
% IF ADVISED OF THE POSSIBILITY OF SUCH DAMAGE.                        %
%                                                                      %
%%%%%%%%%%%%%%%%%%%%%%%%%%%%%%%%%%%%%%%%%%%%%%%%%%%%%%%%%%%%%%%%%%%%%%%%

\addcontentsline{toc}{section}{Stack of Two Graphs}
\psgraphin{appendix/fig/stack2}{15.0}{25.0}{Stack of 2 Graphs}
% \graphin{stack2}{15.0cm}{25.0cm}{Stack of 2 Graphs}
\clearpage
\vfill
\begin{center}
\begin{minipage}[t]{12.0cm}
%%\input{stack2.gle}
{\scriptsize %\footnotesize
\begin{verbatim}
size 18 27  box !  (uncomment box for debugging)    This is STACK2.GLE 
amove 1.3 14.5
begin graph
	size 16 12
	nobox
	xaxis min 0 max 20 dticks 4 dsubticks 1
	xplaces  4 8 12 16 20	! these would be set to match the dates
	xnames "Sep 13" "Sep 23" "Oct 3" "Oct 13" "Oct 23"
	xticks length .2
	yticks length .2
	yaxis min 0 max 75 dticks 10 dsubticks 5 nofirst
	ytitle "Bud burst (%)"
	xtitle ""
	data TEST4.DAT
	d1 lstyle 1 marker wcircle msize .3
	d2 lstyle 1 marker wsquare msize .3
	d3 lstyle 2 marker wcircle msize .3
	d4 lstyle 2 marker wsquare msize .3
end graph
set hei .4 
begin key
	offset 8.2  2.3
	text "Shelter row" lstyle 1 marker circle mscale .6
	text "Shelter row + H " lstyle 1  marker square mscale .6
	text "Middle row" lstyle 2 marker circle mscale .6
	text "Middle row + H" lstyle 2 marker square mscale .6
end key
rmove 12 9 
text (a)

amove 1.3 4 
set font rm hei .5
begin graph
	size 16 12
	nobox
	xaxis min 0 max 20 dticks 4 dsubticks 1
	xplaces  4 8 12 16 20	! these would be set to match the dates
	xnames "Nov 10" "Nov 15"  "Nov 20" "Nov 25"  "Nov 30"
	xticks length .2
	yticks length .2
	yaxis min 0 max 110 dticks 10 dsubticks 5
	ytitle "Flowers per cane at full bloom "
	xtitle ""
	data TEST4.DAT
	d1 lstyle 1 marker circle msize .3
	d2 lstyle 1 marker square msize .3 smooth
	d3 lstyle 2 marker circle msize .3 smooth
	d4 lstyle 2 marker square msize .3
end graph
rmove 12 9 
text (b)

set hei .4 
amove 9 26
text 22.
amove 1.3 3
begin text width 16
	Figure 5. Influence of Hicane on the duration and timing of  (a) bud burst 
	and (b) flowering of kiwifruit.  (Note: this data has been made up)
end text

set just left
amove .5  26
text stack2.gle                                    
\end{verbatim}
} %% \footnotesize
\hfill {\small stack2.gle}\\
\vspace*{-4.0ex} %\vfill
%%\input{test4.dat}
{\footnotesize
\begin{verbatim}
1         10        5         90        12        
3         30        15        80        22        
5         50        25        70        10        
8         80        40        60        15        
10        90        45        55        11        
14        70        35        50        14        
19        80        40        60        6         
\end{verbatim}
}
\hfill {\small test4.dat}\\
\end{minipage}
\end{center}
\vfill
%\psgraphin{stack2}{15.0}{25.0}{Stack of 2 Graphs}
%\graphin{stack2}{15.0cm}{25.0cm}{Stack of 2 Graphs}
%\clearpage

%
% \clearpage
% %%%%%%%%%%%%%%%%%%%%%%%%%%%%%%%%%%%%%%%%%%%%%%%%%%%%%%%%%%%%%%%%%%%%%%%%
%                                                                      %
% GLE - Graphics Layout Engine <http://glx.sourceforge.io/>            %
%                                                                      %
% Modified BSD License                                                 %
%                                                                      %
% Copyright (C) 2009 GLE.                                              %
%                                                                      %
% Redistribution and use in source and binary forms, with or without   %
% modification, are permitted provided that the following conditions   %
% are met:                                                             %
%                                                                      %
%    1. Redistributions of source code must retain the above copyright %
% notice, this list of conditions and the following disclaimer.        %
%                                                                      %
%    2. Redistributions in binary form must reproduce the above        %
% copyright notice, this list of conditions and the following          %
% disclaimer in the documentation and/or other materials provided with %
% the distribution.                                                    %
%                                                                      %
%    3. The name of the author may not be used to endorse or promote   %
% products derived from this software without specific prior written   %
% permission.                                                          %
%                                                                      %
% THIS SOFTWARE IS PROVIDED BY THE AUTHOR "AS IS" AND ANY EXPRESS OR   %
% IMPLIED WARRANTIES, INCLUDING, BUT NOT LIMITED TO, THE IMPLIED       %
% WARRANTIES OF MERCHANTABILITY AND FITNESS FOR A PARTICULAR PURPOSE   %
% ARE DISCLAIMED. IN NO EVENT SHALL THE AUTHOR BE LIABLE FOR ANY       %
% DIRECT, INDIRECT, INCIDENTAL, SPECIAL, EXEMPLARY, OR CONSEQUENTIAL   %
% DAMAGES (INCLUDING, BUT NOT LIMITED TO, PROCUREMENT OF SUBSTITUTE    %
% GOODS OR SERVICES; LOSS OF USE, DATA, OR PROFITS; OR BUSINESS        %
% INTERRUPTION) HOWEVER CAUSED AND ON ANY THEORY OF LIABILITY, WHETHER %
% IN CONTRACT, STRICT LIABILITY, OR TORT (INCLUDING NEGLIGENCE OR      %
% OTHERWISE) ARISING IN ANY WAY OUT OF THE USE OF THIS SOFTWARE, EVEN  %
% IF ADVISED OF THE POSSIBILITY OF SUCH DAMAGE.                        %
%                                                                      %
%%%%%%%%%%%%%%%%%%%%%%%%%%%%%%%%%%%%%%%%%%%%%%%%%%%%%%%%%%%%%%%%%%%%%%%%

%%\psgraphin{stack4b}{15.0}{25.0}{Stack of 4 Graphs}
\graphin{appendix/fig/stack4b}{13.0cm}{25.0cm}{1}
\addcontentsline{toc}{section}{Stack of Four Graphs}
\clearpage
\begin{center}
\begin{minipage}[t]{12.0cm}
%%\input{stack4b.gle}
{\scriptsize %\footnotesize
\begin{verbatim}
size 17 26  ! box  (uncomment box for measuring with ruler) STACK4B.GLE
set font ss
amove 3.5 19
begin graph
	size 10 5
	data TEST.DAT
	title "Top graph, stack4b.gle" hei .5 dist 1
	fullsize
	xaxis min -1 max 6 dticks 1 hei .3 nofirst nolast
	yaxis min 0 max 20 hei .3 dticks 5 
	x2labels on
	d1 marker wsquare  msize .4 lstyle 1
	d2 marker dot  msize 1 lstyle 2
	fill x1,d1 color grid4 xmax 3
	fill x1,d1 color grey20 xmin 3
end graph

rmove 0 -5
begin graph
	size 10 5
	data TEST.DAT
	fullsize
	xaxis min -1 max 6 dticks 1 hei .3 nofirst nolast
	yaxis min 0 max 20 hei .3 dticks 5 nolast
	bar d1,d2 width .2 dist .2 fill grey10,grid3 color black,black
end graph

rmove 0 -5
begin graph
	data TEST.DAT
	fullsize
	size 10 5
	xaxis min -1 max 6 dticks 1 hei .3 nofirst nolast
	yaxis min 0 max 20 hei .3 dticks 5 nolast
	xlabels off
	y2labels on
	d1 marker dot lstyle 2 err 10%
	d2 marker dot lstyle 1 err 10%
end graph

rmove 0 -5
begin graph
	data TEST.DAT
	title "Bottom graph" hei .4 dist 0.0
	fullsize
	size 10 5
	xaxis min -1 max 6 dticks 1 hei .3 nofirst nolast grid
	yaxis min 0 max 20 hei .3 dticks 5 nolast grid
	xticks lstyle 2 lwidth 0.0001  	! sets the grid line style
	yticks lstyle 2 lwidth 0.0001
	d1 marker wsquare  msize .4 lstyle 2 key "Age"
	key hei .5
	d2 marker dot  msize 1 lstyle 1 key "Width"
end graph
\end{verbatim}
}
\hfill {\small stack4b.gle}\\
\end{minipage}
\end{center}
\vfill

%%
% \clearpage
% %%%%%%%%%%%%%%%%%%%%%%%%%%%%%%%%%%%%%%%%%%%%%%%%%%%%%%%%%%%%%%%%%%%%%%%%
%                                                                      %
% GLE - Graphics Layout Engine <http://glx.sourceforge.io/>            %
%                                                                      %
% Modified BSD License                                                 %
%                                                                      %
% Copyright (C) 2009 GLE.                                              %
%                                                                      %
% Redistribution and use in source and binary forms, with or without   %
% modification, are permitted provided that the following conditions   %
% are met:                                                             %
%                                                                      %
%    1. Redistributions of source code must retain the above copyright %
% notice, this list of conditions and the following disclaimer.        %
%                                                                      %
%    2. Redistributions in binary form must reproduce the above        %
% copyright notice, this list of conditions and the following          %
% disclaimer in the documentation and/or other materials provided with %
% the distribution.                                                    %
%                                                                      %
%    3. The name of the author may not be used to endorse or promote   %
% products derived from this software without specific prior written   %
% permission.                                                          %
%                                                                      %
% THIS SOFTWARE IS PROVIDED BY THE AUTHOR "AS IS" AND ANY EXPRESS OR   %
% IMPLIED WARRANTIES, INCLUDING, BUT NOT LIMITED TO, THE IMPLIED       %
% WARRANTIES OF MERCHANTABILITY AND FITNESS FOR A PARTICULAR PURPOSE   %
% ARE DISCLAIMED. IN NO EVENT SHALL THE AUTHOR BE LIABLE FOR ANY       %
% DIRECT, INDIRECT, INCIDENTAL, SPECIAL, EXEMPLARY, OR CONSEQUENTIAL   %
% DAMAGES (INCLUDING, BUT NOT LIMITED TO, PROCUREMENT OF SUBSTITUTE    %
% GOODS OR SERVICES; LOSS OF USE, DATA, OR PROFITS; OR BUSINESS        %
% INTERRUPTION) HOWEVER CAUSED AND ON ANY THEORY OF LIABILITY, WHETHER %
% IN CONTRACT, STRICT LIABILITY, OR TORT (INCLUDING NEGLIGENCE OR      %
% OTHERWISE) ARISING IN ANY WAY OUT OF THE USE OF THIS SOFTWARE, EVEN  %
% IF ADVISED OF THE POSSIBILITY OF SUCH DAMAGE.                        %
%                                                                      %
%%%%%%%%%%%%%%%%%%%%%%%%%%%%%%%%%%%%%%%%%%%%%%%%%%%%%%%%%%%%%%%%%%%%%%%%

%%\psgraphin{stack1}{15.0}{23.0}{Shadow Graph}
\graphin{appendix/fig/stack1}{15.0cm}{23.0cm}{1}
\addcontentsline{toc}{section}{Shadow Graph}
\clearpage
\begin{center}
\begin{minipage}[t]{12.0cm}
%%\input{stack1.gle}
{\scriptsize %\footnotesize
\begin{verbatim}
size 17 23 	! This is STACK1.GLE,  a single graph, with shading background
set hei .4
amove .3 .3
text stack1.gle

amove 1.5 5
box 15 16 fill grey60
rmove -1 1
box 15 16 fill white
rmove 2 4
box 11 8 fill grey5
set font ss
begin graph
	fullsize
	size 11 8
	title "Data transfer on 9600 bit/sec line"
	xtitle "Time (seconds)"
	ytitle "K Bits"
	data test.dat
	d1 line marker wsquare mscale 2
	xaxis min -1 max 6
	yaxis min 0 max 11
end graph

\end{verbatim}
}
\hfill {\small stack1.gle}\\
\end{minipage}

\end{center}
\vfill

%%
% \clearpage
% %%%%%%%%%%%%%%%%%%%%%%%%%%%%%%%%%%%%%%%%%%%%%%%%%%%%%%%%%%%%%%%%%%%%%%%%
%                                                                      %
% GLE - Graphics Layout Engine <http://glx.sourceforge.io/>            %
%                                                                      %
% Modified BSD License                                                 %
%                                                                      %
% Copyright (C) 2009 GLE.                                              %
%                                                                      %
% Redistribution and use in source and binary forms, with or without   %
% modification, are permitted provided that the following conditions   %
% are met:                                                             %
%                                                                      %
%    1. Redistributions of source code must retain the above copyright %
% notice, this list of conditions and the following disclaimer.        %
%                                                                      %
%    2. Redistributions in binary form must reproduce the above        %
% copyright notice, this list of conditions and the following          %
% disclaimer in the documentation and/or other materials provided with %
% the distribution.                                                    %
%                                                                      %
%    3. The name of the author may not be used to endorse or promote   %
% products derived from this software without specific prior written   %
% permission.                                                          %
%                                                                      %
% THIS SOFTWARE IS PROVIDED BY THE AUTHOR "AS IS" AND ANY EXPRESS OR   %
% IMPLIED WARRANTIES, INCLUDING, BUT NOT LIMITED TO, THE IMPLIED       %
% WARRANTIES OF MERCHANTABILITY AND FITNESS FOR A PARTICULAR PURPOSE   %
% ARE DISCLAIMED. IN NO EVENT SHALL THE AUTHOR BE LIABLE FOR ANY       %
% DIRECT, INDIRECT, INCIDENTAL, SPECIAL, EXEMPLARY, OR CONSEQUENTIAL   %
% DAMAGES (INCLUDING, BUT NOT LIMITED TO, PROCUREMENT OF SUBSTITUTE    %
% GOODS OR SERVICES; LOSS OF USE, DATA, OR PROFITS; OR BUSINESS        %
% INTERRUPTION) HOWEVER CAUSED AND ON ANY THEORY OF LIABILITY, WHETHER %
% IN CONTRACT, STRICT LIABILITY, OR TORT (INCLUDING NEGLIGENCE OR      %
% OTHERWISE) ARISING IN ANY WAY OUT OF THE USE OF THIS SOFTWARE, EVEN  %
% IF ADVISED OF THE POSSIBILITY OF SUCH DAMAGE.                        %
%                                                                      %
%%%%%%%%%%%%%%%%%%%%%%%%%%%%%%%%%%%%%%%%%%%%%%%%%%%%%%%%%%%%%%%%%%%%%%%%

%%\psgraphin{loop}{15.0}{25.0}{Loop, Subroutines}
\graphin{appendix/fig/loop}{15.0cm}{25.0cm}{1}
\addcontentsline{toc}{section}{Loop, Subroutines}
\clearpage
\begin{center}
\begin{minipage}[t]{12.0cm}
%%\input{loop.gle}
{\scriptsize %\footnotesize
\begin{verbatim}
size 17 17 box !18 24 box		! This is LOOP.GLE, shows programming examples
amove 5 14
set just center hei 1
text Subroutines, Loops
set just left hei .7
set hei .9

! First define a single neuron
sub neuron xx yy
   amove xx yy 
   begin scale .5 .5
      curve  0  .4  .4  .4  .4 -.4  .4  .4  .4 -.8 -.4 -.8 -.4 .4
      curve -.4 -.8 -.4  .4 -.4 -.4 -.4  .4  0  .4   
      closepath
   end scale
end sub

! Draw the neurons
for xx = 0 to 1
   for i = 1 to 6
      @neuron xx*6+2 i*2+1.2
   next i
next xx

! Draw the lines between them
for i = 1 to 6
   for j = 1 to 6
      amove 3 i*2+1
      aline 7.5 j*2+1
   next j
next i

!--------------------------------------------------------------------------
include ziptext.gle
amove 12 1
@ziptext "Hello"
amove 13 4
begin rotate 20
   @ziptext "Green trees"
end rotate
\end{verbatim}
}
\hfill {\small loop.gle}\\
%
%%\input{loop.gle}
{\scriptsize %\footnotesize
\begin{verbatim}
sub ziptext tt$
gsave
for i = 1 to 0 step -.05
	set color (i)
	write tt$ 
	rmove -.05 .025
next i
set color white
write tt$ 
grestore
end sub

\end{verbatim}
}
\hfill {\small ziptext.gle}\\
%
\end{minipage}
\end{center}
\vfill

%%
% \clearpage
% %%%%%%%%%%%%%%%%%%%%%%%%%%%%%%%%%%%%%%%%%%%%%%%%%%%%%%%%%%%%%%%%%%%%%%%%
%                                                                      %
% GLE - Graphics Layout Engine <http://glx.sourceforge.io/>            %
%                                                                      %
% Modified BSD License                                                 %
%                                                                      %
% Copyright (C) 2009 GLE.                                              %
%                                                                      %
% Redistribution and use in source and binary forms, with or without   %
% modification, are permitted provided that the following conditions   %
% are met:                                                             %
%                                                                      %
%    1. Redistributions of source code must retain the above copyright %
% notice, this list of conditions and the following disclaimer.        %
%                                                                      %
%    2. Redistributions in binary form must reproduce the above        %
% copyright notice, this list of conditions and the following          %
% disclaimer in the documentation and/or other materials provided with %
% the distribution.                                                    %
%                                                                      %
%    3. The name of the author may not be used to endorse or promote   %
% products derived from this software without specific prior written   %
% permission.                                                          %
%                                                                      %
% THIS SOFTWARE IS PROVIDED BY THE AUTHOR "AS IS" AND ANY EXPRESS OR   %
% IMPLIED WARRANTIES, INCLUDING, BUT NOT LIMITED TO, THE IMPLIED       %
% WARRANTIES OF MERCHANTABILITY AND FITNESS FOR A PARTICULAR PURPOSE   %
% ARE DISCLAIMED. IN NO EVENT SHALL THE AUTHOR BE LIABLE FOR ANY       %
% DIRECT, INDIRECT, INCIDENTAL, SPECIAL, EXEMPLARY, OR CONSEQUENTIAL   %
% DAMAGES (INCLUDING, BUT NOT LIMITED TO, PROCUREMENT OF SUBSTITUTE    %
% GOODS OR SERVICES; LOSS OF USE, DATA, OR PROFITS; OR BUSINESS        %
% INTERRUPTION) HOWEVER CAUSED AND ON ANY THEORY OF LIABILITY, WHETHER %
% IN CONTRACT, STRICT LIABILITY, OR TORT (INCLUDING NEGLIGENCE OR      %
% OTHERWISE) ARISING IN ANY WAY OUT OF THE USE OF THIS SOFTWARE, EVEN  %
% IF ADVISED OF THE POSSIBILITY OF SUCH DAMAGE.                        %
%                                                                      %
%%%%%%%%%%%%%%%%%%%%%%%%%%%%%%%%%%%%%%%%%%%%%%%%%%%%%%%%%%%%%%%%%%%%%%%%

%%\psgraphin{stack1}{15.0}{23.0}{Shadow Graph}
\graphin{appendix/fig/join}{14.0cm}{19.0cm}{1}
\addcontentsline{toc}{section}{Joining Objects}
\clearpage
\begin{center}
\begin{minipage}[t]{12.0cm}
%%\input{join.gle}
{\scriptsize %\footnotesize
\begin{verbatim}
size 18 24 box	! This is JOIN.GLE, demonstrates the join commands.
set hei .3
amove .3 .3 
text join.gle

set hei 1.3 just center
amove 9 21
text Joining named objects

set hei .8 lwidth .1
amove  8.8  18.8 
begin box name line
   rline  0 -8.2 
end box
set lwidth 0

amove  9.67745  8.7603 
begin box   name cir   nobox
   circle 1
end box
circle 1
amove  12.7461  4.59445 
begin box  name hi   add .2
   text Hi there
end box
amove  3.14401  17.0566 
begin box name grv add .2
   text GRV
end box
amove  14.9533  12.2986 
begin box name chv add .2
   text CHV
end box
amove  12.5179  17.6103 
begin box name cheese add .2
   text Cheese
end box
amove  2.7381  11.229 
begin box name goats add .2
   text Goats
end box
! Now draw the lines between objects
join chv - goats
join grv -> line.h          ! Note: ".h" means to join horizontally.
join line.h <-> cheese.tl
join cheese.tr - chv.tr
join cir.ci <- hi           ! Note: ".ci" is used for circles
join cir.ci <- chv
join cir.ci <- goats
!-----------------------------------------------------------------------------
! A pie slice   (Will only work on PostScript it uses BEGIN PATH FILL).
sub pie ang1 ang2 radius color$
begin path fill color$ stroke
   rmove 0 0
   arc radius ang1 ang2 
   closepath
end path
end sub

amove 3 2
@pie 0 10  2 "grey10"
@pie 10 40 2 "grey20"
@pie 40 120 2 "blue"
\end{verbatim}
}
\hfill {\small join.gle}\\
\end{minipage}

\end{center}
\vfill

%%
% \clearpage
% %%%%%%%%%%%%%%%%%%%%%%%%%%%%%%%%%%%%%%%%%%%%%%%%%%%%%%%%%%%%%%%%%%%%%%%%
%                                                                      %
% GLE - Graphics Layout Engine <http://glx.sourceforge.io/>            %
%                                                                      %
% Modified BSD License                                                 %
%                                                                      %
% Copyright (C) 2009 GLE.                                              %
%                                                                      %
% Redistribution and use in source and binary forms, with or without   %
% modification, are permitted provided that the following conditions   %
% are met:                                                             %
%                                                                      %
%    1. Redistributions of source code must retain the above copyright %
% notice, this list of conditions and the following disclaimer.        %
%                                                                      %
%    2. Redistributions in binary form must reproduce the above        %
% copyright notice, this list of conditions and the following          %
% disclaimer in the documentation and/or other materials provided with %
% the distribution.                                                    %
%                                                                      %
%    3. The name of the author may not be used to endorse or promote   %
% products derived from this software without specific prior written   %
% permission.                                                          %
%                                                                      %
% THIS SOFTWARE IS PROVIDED BY THE AUTHOR "AS IS" AND ANY EXPRESS OR   %
% IMPLIED WARRANTIES, INCLUDING, BUT NOT LIMITED TO, THE IMPLIED       %
% WARRANTIES OF MERCHANTABILITY AND FITNESS FOR A PARTICULAR PURPOSE   %
% ARE DISCLAIMED. IN NO EVENT SHALL THE AUTHOR BE LIABLE FOR ANY       %
% DIRECT, INDIRECT, INCIDENTAL, SPECIAL, EXEMPLARY, OR CONSEQUENTIAL   %
% DAMAGES (INCLUDING, BUT NOT LIMITED TO, PROCUREMENT OF SUBSTITUTE    %
% GOODS OR SERVICES; LOSS OF USE, DATA, OR PROFITS; OR BUSINESS        %
% INTERRUPTION) HOWEVER CAUSED AND ON ANY THEORY OF LIABILITY, WHETHER %
% IN CONTRACT, STRICT LIABILITY, OR TORT (INCLUDING NEGLIGENCE OR      %
% OTHERWISE) ARISING IN ANY WAY OUT OF THE USE OF THIS SOFTWARE, EVEN  %
% IF ADVISED OF THE POSSIBILITY OF SUCH DAMAGE.                        %
%                                                                      %
%%%%%%%%%%%%%%%%%%%%%%%%%%%%%%%%%%%%%%%%%%%%%%%%%%%%%%%%%%%%%%%%%%%%%%%%

%%\psgraphin{stack1}{15.0}{23.0}{Shadow Graph}
\graphin{appendix/fig/clip}{14.0cm}{19.0cm}{1}
\addcontentsline{toc}{section}{Clipping}
\clearpage
\begin{center}
\begin{minipage}[t]{12.0cm}
%%\input{clip.gle}
{\scriptsize %\footnotesize
\begin{verbatim}
size 18 24	! This is CLIP.GLE, sets up an arbitrary clipping region.
set hei .3	! NOTE: CLIPPING ONLY WORKS ON POSTSCRIPT DEVICES.
amove 2.3 .3	! So this won't work on the pc-screen.
text clip.gle	

set hei 10
amove 2 2 
begin clip	! Save default clipping region
   begin path clip	! Set up the clipping region
      text GLE
   end path
   gsave
      amove 2 2 	
      text www	! Draw some text which will be clipped
      for i =	1 to 90	! Draw some lines which will be clipped
         amove 2	2 
         rotate -1
         rline 0 30
      next i
   grestore
end clip ! restore default clipping region
amove 0 0 
rline 10 10 	! draw a line which won't be clipped.
!-----------------------------------------------------------------------------
!Here's a simple example of how to set up a clipping region which consists
!of two squares. Try this out
amove 1 11
begin origin
   begin path clip stroke
      amove 2 2 
      box 4 4 
      amove 7 2 
      box 4 4 
   end path
   amove 2 2 
   set hei 4
   text Here is clipped text
end origin
\end{verbatim}
}
\hfill {\small clip.gle}\\
\end{minipage}

\end{center}
\vfill

%%
%%\cleardoublepage

